\documentclass{article}
\usepackage[utf8]{inputenc}

\title{Linear Regression Notes}
\author{Conor Falvey}

\usepackage{natbib}
\usepackage{graphicx}

\begin{document}

\maketitle

\section{Thursday August 22}
\textbf{Unsupervised learning:} Analyzes structures of X. 
Supervised Learning: Also have a label or response for every observation.
$y = (y_1, y_2, \ldots, y_n)$\\
\textbf{Goal:} Identify a function $f$ that relates $y$ to $x$: $y = f(x) + \epsilon$
The choice of $f(x)$ depends on several objectives:
\begin{itemize}
    \item Explanatory or Predictive performance?
    \begin{itemize}
        \item Trade-off on outputs
    \end{itemize}
    \item Parametric or non-parametric?
    \begin{itemize}
        \item Parametric: Characterized by parameters, Summary of a population
        \item Non-parametric: No parameters. (Technically infinite)
    \end{itemize}
    \item Functional or Statistical Relationship?
    \begin{itemize}
        \item Functional Relationship: $y = mx + b$. Function describes a relationship. It has a deterministic output.
        \item Statistical Relationship: $y = mx + b + \epsilon$. Involves randomness.
        \begin{itemize}
            \item $\epsilon$ is our uncertainty in the relationship we chose.
            \item $\epsilon$ is often referred to as error.
        \end{itemize}
    \end{itemize}
\end{itemize}
\textbf{Overall Question:} How do we choose $f(x)$?
The key to data science and modeling is answering this question with data in a 'smart' way. In this class we'll focus on 'linear' families of models. Why linear?
\begin{itemize}
    \item Easy to understand
    \begin{itemize}
        \item Highly interpretable
        \item Explanatory inferences are straightforward
    \end{itemize}
    \item Lines can be used to approximate more complex relationships (Think of Taylor expansions.)
\end{itemize}
When choosing a model $f$, we are usually concerned with one of two primary goals: \textbf{prediction} or \textbf{inference}. It is possible to choose a model that is reasonably well-calibrated for both prediction and inference.\\

\textbf{Prediction:}\\
\textbf{Main Objective:} Predicting new $Y$ using $\hat{Y} = \hat{f}(X)$\\
\textbf{Model choice:} Models that have the highest prediction performance. Often 'black-box methods, which have no concern with the exact form of $\hat{f}$\\

\textbf{Inference:}\\
\textbf{Main Objective:} Understand the relationship between $X$ and $Y$:
\begin{itemize}
    \item Variable Selection: What predictors are most associated with the response?
    \item Focus on functional relationships of $Y$ and $X$
    \item Strive for parsimony. KISS: Keep It Simple Stupid
\end{itemize}
Model choice: Models that have high interpretability. Often parametric methods, which explicitly dictate the form of $\hat{f}$ via parameters.\\


\textbf{Parametric Methods:}
\begin{itemize}
    \item Makes an assumption about the functional form of $f$.
    \item Identifying $f$ reduces the estimation of a set of parameters.
    \item Often simpler than estimating the entire function.
    \item Caution! Can result in over-simplistic models.
    \item E.g. linear regression, logistic regression\\
\end{itemize}

\textbf{Non-parametric Methods:}
\begin{itemize}
    \item Not restricted to assumptions about the functional form of $f$.
    \item Can accurately fit a wider range of possible shapes/forms of $f$.
    \item Often requires a very large number of observations.
    \item Caution: Can overfit data quickly.
    \item E.g. Polynomial splines, smoothing splines\\
\end{itemize}
\textbf{Question:}
How well does your method perform on new data, i.e. data you have not seen during learning?
E.g. You get a new data instance
$X_{new} = ($no history of cancer, smoker, male$)$
Can you asses the goodness of your throat cancer predictor?\\
\textbf{Past Observed Data:}\\
$(x_{1}, y_{1}), (x_{2}, y_{2},\ldots)$\\
Previous Data
$X = ($History of cancer, smoker, sex$)$\\
$y = \{$0 if no cancer, 1 if cancer$\}$\\
\textbf{New Data:}\\
Only with features: $X_{new}$
Without having $y_{new}$, how can we assess the performance of your prediction $\hat{y}_{new}$?\\
\textbf{Options:}
\begin{enumerate}
    \item If we could get a true diagnosis, $y_{new}$, then measure the difference $(y_{new} - \hat{y}_{new})^2$.
    \item If we can't get $y_{new}$, then use cross-validation from past data. Go back to past data and use it to assess current performance.
\end{enumerate}
\textbf{Training vs. Test Set Validation}
\begin{enumerate}
    \item Divide the data $(X, Y)$ into two sets:
    \begin{itemize}
        \item Training Set: $(X_{train}, Y_{train})$
        \item Test set: $(X_{test}, Y_{test})$
    \end{itemize}
    \item Use training set to produce a predictor $\hat{f}()$ via $y_{train} = f(X_{train}) + \epsilon$. Learn/train a model on training set $\hat{f}$.
    \item Use the test set to evaluate the performance of predictor: $\hat{y}_{test} = \hat{f}(X_{test})$. Access the difference between $\hat{y}_{test}$ and $y_{test}$.
\end{enumerate}
\textbf{Choosing Training and Test Set}
\begin{enumerate}
    \item Random Sampling: Choose a test set at random from data and test all models on the same set.
    \item k-fold cross-validation: Split data into k subsets. In turn, treat each subset as held-out and train on the remaining. Performance is evaluated as the average performance of each k test sets. (Compute difference across each iteration and average over all iterations.)
    \item Leave-one-out cross-validation: Special case of k-fold cross validation where $k = n$ and test sets are of size 1.
\end{enumerate}
\textbf{Important considerations:}
Never let information from the test set make its way into the training data! This is the \#1 most common mistake in model assessment.\\

Difference between \textbf{prediction and estimation error}:
\begin{enumerate}
    \item Prediction error: Error associated with predictions on the test set $y_{test} = \hat{y}_{test}$
    \item Estimation error: Error associated with estimates in the training set $y_{train} = \hat{y}_{train}$
    \item By Splitting the data in cross-validation, the variance of the estimated regression coefficients can increase if the data set is not large
    \item Once a model has been validated and compared against other potential models, we typically use the entire data set for estimating the final regression model
    \item Sometimes the training and test sets are chosen in a systematic way (e.g. up-sampling and down-sampling) so as to avoid bias in the analysis.
\end{enumerate}
When running cross-validation, there are \textbf{2 types of error:}
\begin{enumerate}
    \item Training error aka estimation error $(y_{train} - \hat{y}_{train})^2$. The error associated with estimation/interpolation.
    \item Test error $(y_{test} - \hat{y}_{test})^2$. The error associated with prediction and extrapolation.
\end{enumerate}
\textbf{Mean Squared Error:}\\
\begin{center}
    MSE$(f) = \frac{1}{n}\sum^{n}_{i = 1}(y_{i} - f(x_{i}))^2$\\
\end{center} You can use any positive even exponent, we just want all summation parts to end positive. The MSE measure acts as a yardstick for model assessment for continuous data.\\
Note: In prediction, the mean squared difference between $y_{new}$ and $\hat{f}(x_{new})$ is known as the \textbf{mean square prediction error (MSPE)}.\\
Note: There are many choices for measuring accuracy. The choice depends on the possible values of $y$.\\
\textbf{Assessing Model Accuracy:}\\
Let $\Omega$ = index (which rows of X) that represent the training set.\\
Let $\Theta = \Omega^{c}$ = index that represent the test set.\\
\textbf{General Approach:}\\
Estimate model $\hat{f}$:
\begin{center}
$\hat{f} = argmin_{f}(\frac{1}{|\Omega|}\sum_{j \in \Omega}(y_{j} - f(x_{j}))^2) = argmin_{f}($MSE$(f))$
\end{center}
Evaluate model on test set:
\begin{center}
MSPE$(\hat{f}) = \frac{1}{|\Theta|}sum_{j \in \Theta}(y_{j} - \hat{f}(x_{j}))^2$\\
\end{center}
\textbf{MSE vs. MSPE: Noticeable Trends}
\begin{itemize}
    \item Fact: $E[$MSPE$(f(X_{test})] \geq E[$MSE$(f(X_{train}))]]$
    \item Trend 1: After a certain point in complexity, there is an inverse relationship between MSE and MSPE. We say that a model is overfitting the data when we are in this range of complexity.
    \item Trend 2: The MSE tends to decrease as complexity increases.
\end{itemize}
\textbf{Bias-Variance Trade-Off:}
\begin{itemize}
    \item An important means of understanding \begin{center}
        MSPE$(\hat{f}) = \frac{1}{|\Theta|}sum_{j \in \Theta}(y_{j} - \hat{f}(x_{j}))^2$
    \end{center} 
    comes from the following decomposition for new data $(X_{0}, y_{0})$: 
    \begin{center}
        $E[$MSPE$(\hat{f})] = E[(y_{0} - \hat{f}(X_{0})]^2 + Var(\hat{f}(X_{0})) + Var(\epsilon)$ = $Bias(\hat{f}(X_{0}))^2 + Var(\hat{f}(X_{0})) + Var(\epsilon)$
    \end{center}
    \item Result: the expected MSPE of a model is a function of the bias and variance of $\hat{f}$, as well as the variance of the error term $\epsilon$.
\end{itemize}
\textbf{Components of MSPE:}
\begin{enumerate}
    \item $Bias(\hat{f}(X_{0}))^2$: quantifies the distance between the model and the truth. Non-negative. Generally decreases as the model becomes more complex.
    \item $Var(\hat{f}(X_{0}))$: quantifies the variance of the model. Non-negative. High variance implies that the model is highly sensitive to small changes in the training data. Generally increases as the model becomes more complex.
    \item $Var(\epsilon)$: variance of the error terms. Non-negative. Is not affected by the complexity of the model, constant.
\end{enumerate}
\textbf{Bias-Variance Trade-Off:}\\
Resulting Trade-Off: Since $Var(\epsilon)$ is constant, we'd like to choose a model with minimum bias and minimum variance.\\
Primary Issues:
\begin{enumerate}
    \item Both the bias and variance terms are non-negative
    \item The bias and variance terms are often inversely related
    \item The bias and variance change at different rates
\end{enumerate}
\textbf{Solution:}\\
Decide what is important in the application (prediction vs. interpretation) and choose model accordingly. Seek optimal model complexity of possible.\\

\textbf{Regression Model Function:}
\begin{itemize}
    \item Simple Linear Regression
    \begin{itemize}
        \item Simple because we consider the relationship between a single feature, $x_{1}$, and response, $y_{1}$.
        \begin{itemize}
            \item Model $y_{i} = \beta_{0} + \beta_{1}x_{i} + \epsilon_{i}$
            \item Data: Pairs $(x_{1}, y_{1}), (x_{2}, y_{2}), \ldots, (x_{n}, y_{n})$
            \item Assumptions (1):
            \begin{itemize}
                \item $E[\epsilon_{i}] = 0$
                \item $Var(\epsilon_{i}) = \sigma$
                \item $Cov(\epsilon_{i}, \epsilon_{j}) = 0$
            \end{itemize}
            \item Assumptions (2):
            \begin{itemize}
                \item $E_{i} \stackrel{IID}{\sim} N(0, \sigma^2)$. $E_{i}$ = model error on observation i.
            \end{itemize}
            \item Components of Simple Linear Regression Model:
            \begin{itemize}
                \item $y_{1}, \ldots, y_{n}$ are continuous responses
                \item $x_{1}, \ldots, x_{n}$ are feature values
                \item $\beta_{0}$ is the intercept of the model
                \item $\beta_{2}$ is the effect on y from a single unit change in $x$
                \item $\epsilon$ is the error term
                \item $\sigma^2$ is the variance of the error
                \item $\beta_{0}, \beta_{1}, \sigma^2$ are unknown parameters of the model
            \end{itemize}
        \end{itemize}
    \end{itemize}
\end{itemize}

\end{document}
